\input{../tex-inputs/latex-korrekturansicht-vorspann}

\section[Berta Zuckerkandl an Arthur Schnitzler, {[}7. 11. 1911?{]}]{L04004 Berta Zuckerkandl an Arthur Schnitzler, {[}7. 11. 1911?{]}}
\nopagebreak\mylabel{L04004v}
\rehead{ }\normalsize\beginnumbering\briefempfaengerindex{Schnitzler, Arthur@\textsc{Schnitzler, Arthur}!zzzZuckerkandl, Berta@\emph{von Berta Zuckerkandl}!1911-11-072@{{[}7. 11. 1911?{]}}|(be}
\toendnotes[C]{\smallbreak\pagebreak[2]}
\correspDesc{Versand  durch Berta Zuckerkandl am [7. 11. 1911?] in Wien
\newline{}Erhalt  durch Arthur Schnitzler im Zeitraum [7. 11. 1911 – 10. 11. 1911?] in Wien}\toendnotes[C]{\smallbreak}
\Standort{CUL, Schnitzler, B 200.}
\physDesc{Brief, 1 Blatt, 4 Seiten, 686 Zeichen
\newline{}Handschrift: schwarze Tinte, lateinische Kurrent
\newline{}Beilage: beschnittener Zeitungsausschnitt aus \textcolor{green}{Le Figaro}\pwindex{Le Figaro@\emph{Le Figaro}|pw} mit ungezeichneter Notiz von \textcolor{blue}{Berta Zuckerkandl}\pwindex{Zuckerkandl, Berta 13.\,4.\,1864 Wien – 16.\,10.\,1945 Paris@\textsc{Zuckerkandl, Berta} (13.\,4.\,1864 Wien – 16.\,10.\,1945 Paris), \emph{Schriftstellerin, Journalistin, Übersetzerin}|pw}: \emph{\textcolor{green}{Courrier des Théatres. Au jour
                                       le jour. [Shakespeare disait…]}\pwindex{Zuckerkandl, Berta 13.\,4.\,1864 Wien – 16.\,10.\,1945 Paris@\textsc{Zuckerkandl, Berta} (13.\,4.\,1864 Wien – 16.\,10.\,1945 Paris), \emph{Schriftstellerin, Journalistin, Übersetzerin}!Courrier des Théatres. Au jour le jour. Shakespeare disait [...] [Über die Wiener Uraufführung von Das weite Land]@\strich\emph{Courrier des Théatres. Au jour le jour. Shakespeare disait [...] [Über die Wiener Uraufführung von Das weite Land]}|pwk}}. In: \emph{\textcolor{green}{Le Figaro}\pwindex{Le Figaro@\emph{Le Figaro}|pwk}}, Jg. 57, Serie 3, Nr. 304,
                                       31. 10. 1911, S. 6. 
\newline{}Schnitzler: mit Bleistift Vermerk: »\textsc{Zucker}{[}kandl{]}« }\toendnotes[C]{\smallbreak}
\pstart
           \raggedleft{}{\pb}\label{K_L04004-1v}\edtext{Dienstag}{\lemma{\textnormal{\emph{Dienstag}}}\Cendnote{\textnormal{Der vorliegende Brief ist nicht genauer
                     datiert. Die beigelegte \textcolor{green}{Zeitungsnotiz}\pwindex{Zuckerkandl, Berta 13.\,4.\,1864 Wien – 16.\,10.\,1945 Paris@\textsc{Zuckerkandl, Berta} (13.\,4.\,1864 Wien – 16.\,10.\,1945 Paris), \emph{Schriftstellerin, Journalistin, Übersetzerin}!Courrier des Théatres. Au jour le jour. Shakespeare disait [...] [Über die Wiener Uraufführung von Das weite Land]@\strich\emph{Courrier des Théatres. Au jour le jour. Shakespeare disait [...] [Über die Wiener Uraufführung von Das weite Land]}|pwkv} erschien
                     am 31. 10. 1911, die dadurch angestoßenen Anfragen um Übersetzung von \emph{\textcolor{green}{Das weite Land}\pwindex{Schnitzler, Arthur 15. 5. 1862 Wien – 21. 10. 1931 ebd.@\textsc{Schnitzler, Arthur} (15. 5. 1862 Wien – 21. 10. 1931 ebd.), \emph{Schriftsteller, Mediziner}!weite Land. Tragikomödie in fünf Akten@\strich\emph{Das weite Land. Tragikomödie in fünf Akten}|pwk}} stammen vom 31. 10. 1911 (\textcolor{blue}{Wilhelm Bauer}\pwindex{Bauer, Wilhelm 27.\,11.\,1854 Zollingen – 11.\,9.\,1923 Paris@\textsc{Bauer, Wilhelm} (27.\,11.\,1854 Zollingen – 11.\,9.\,1923 Paris)|pwk})
                     und 1. 11. 1911
                     (\textcolor{blue}{Maurice Rémon}\pwindex{Rémon, Maurice 27.\,11.\,1861 Paris – 20.\,6.\,1945 Mérignac@\textsc{Rémon, Maurice} (27.\,11.\,1861 Paris – 20.\,6.\,1945 Mérignac), \emph{Übersetzer}|pwk}), so dass der darauffolgende Dienstag 
                     den mutmaßlichen Tag der Abfassung dieses Schreibens darstellt. \textcolor{blue}{Schnitzler} hatte \textcolor{pink}{Wien}\oindex{Wien@\textbf{Wien}, \emph{Verwaltungsgebiet}|pwk} am
                        30. 10. 1911
                     Richtung \textcolor{pink}{Prag}\oindex{Prag@\textbf{Prag}, \emph{Land}|pwk} verlassen und kehrte erst am
                        17. 11. 1911
                     nach Stationen in \textcolor{pink}{Dresden}\oindex{Dresden@\textbf{Dresden}|pwk}, \textcolor{pink}{Berlin}\oindex{Berlin@\textbf{Berlin}, \emph{Hauptstadt}|pwk}, \textcolor{pink}{Hamburg}\oindex{Hamburg@\textbf{Hamburg}|pwk}, \textcolor{pink}{München}\oindex{München@\textbf{München}|pwk} und \textcolor{pink}{Garmisch-Partenkirchen}\oindex{Garmisch-Partenkirchen@\textbf{Garmisch-Partenkirchen}, \emph{Hauptstadt}|pwk} nach Hause zurück. Am 18. 11. 1911 traf er
                        \textcolor{blue}{Berta Zuckerkandl}\pwindex{Zuckerkandl, Berta 13.\,4.\,1864 Wien – 16.\,10.\,1945 Paris@\textsc{Zuckerkandl, Berta} (13.\,4.\,1864 Wien – 16.\,10.\,1945 Paris), \emph{Schriftstellerin, Journalistin, Übersetzerin}|pwk} und besprach die
                     Angelegenheit. In dieser Zeitspanne lag ein zweiter Dienstag – der 14. 11. 1911 – der ebenfalls für
                     die Datierung des Schreibens in Frage kommt.}}}\label{K_L04004-1}.\pend
           
\pstart{}Verehrter Herr Doktor!\pend\vspace{0.5em}
\pstart
           Anbei sende ich Ihnen \label{K_L04004-2v}\edtext{die \textcolor{green}{Notiz}\pwindex{Zuckerkandl, Berta 13.\,4.\,1864 Wien – 16.\,10.\,1945 Paris@\textsc{Zuckerkandl, Berta} (13.\,4.\,1864 Wien – 16.\,10.\,1945 Paris), \emph{Schriftstellerin, Journalistin, Übersetzerin}!Courrier des Théatres. Au jour le jour. Shakespeare disait [...] [Über die Wiener Uraufführung von Das weite Land]@\strich\emph{Courrier des Théatres. Au jour le jour. Shakespeare disait [...] [Über die Wiener Uraufführung von Das weite Land]}|pwv}{}\ledrightnote{{$\rightarrow$}\emph{\textcolor{green}{Courrier des Théatres. Au jour le jour. Shakespeare disait [...] [Über die Wiener Uraufführung von Das weite Land]}}}}{\lemma{\textnormal{\emph{die Notiz}}}\Cendnote{\textnormal{\emph{\textcolor{green}{Courrier des Théatres. Au jour le jour.
                        Shakespeare disait […]}\pwindex{Zuckerkandl, Berta 13.\,4.\,1864 Wien – 16.\,10.\,1945 Paris@\textsc{Zuckerkandl, Berta} (13.\,4.\,1864 Wien – 16.\,10.\,1945 Paris), \emph{Schriftstellerin, Journalistin, Übersetzerin}!Courrier des Théatres. Au jour le jour. Shakespeare disait [...] [Über die Wiener Uraufführung von Das weite Land]@\strich\emph{Courrier des Théatres. Au jour le jour. Shakespeare disait [...] [Über die Wiener Uraufführung von Das weite Land]}|pwk}}. In: \emph{\textcolor{green}{Le
                        Figaro}\pwindex{Le Figaro@\emph{Le Figaro}|pwk}}, Jg. 57, Serie 3, Nr. 304, 31. 10. 1911,
                     S. 6, siehe unten.}}}\label{K_L04004-2} die ich im »\textcolor{green}{Figaro}\pwindex{Le Figaro@\emph{Le Figaro}|pw}{}\ledrightnote{\textcolor{green}{Le Figaro}}« erscheinen liess. Daraufhin haben Sie die mir \label{K_L04004-3v}\edtext{eingesendeten Briefe}{\lemma{\textnormal{\emph{eingesendeten Briefe}}}\Cendnote{\textnormal{\textcolor{blue}{Wilhelm Bauer}\pwindex{Bauer, Wilhelm 27.\,11.\,1854 Zollingen – 11.\,9.\,1923 Paris@\textsc{Bauer, Wilhelm} (27.\,11.\,1854 Zollingen – 11.\,9.\,1923 Paris)|pwk} an \textcolor{blue}{Arthur Schnitzler}, 31. 10. 1911, und \textcolor{blue}{Maurice Rémon}\pwindex{Rémon, Maurice 27.\,11.\,1861 Paris – 20.\,6.\,1945 Mérignac@\textsc{Rémon, Maurice} (27.\,11.\,1861 Paris – 20.\,6.\,1945 Mérignac), \emph{Übersetzer}|pwk} an \textcolor{blue}{Arthur
                     Schnitzler}, 1. 11. 1911, vgl. Karl Zieger: \emph{Arthur
                        Schnitzler et la France 1894–1938. Enquête sur une réception},
                     Villeneuve d’Ascq: \emph{Presses Universitaires du
                        Septentrion} 2012,
                  S. 190.}}}\label{K_L04004-3}{ }{\pb}wol erhalten. Besonders Herrn \textcolor{blue}{Rémons}\pwindex{Rémon, Maurice 27.\,11.\,1861 Paris – 20.\,6.\,1945 Mérignac@\textsc{Rémon, Maurice} (27.\,11.\,1861 Paris – 20.\,6.\,1945 Mérignac), \emph{Übersetzer}|pw}{}\ledrightnote{\textcolor{blue}{Maurice Rémon}}{ }\uline{Hinweis} auf \textcolor{blue}{Guitry}\pwindex{Guitry, Lucien 13.\,12.\,1860 Paris – 1.\,6.\,1925 ebd.@\textsc{Guitry, Lucien} (13.\,12.\,1860 Paris – 1.\,6.\,1925 ebd.), \emph{Schriftsteller, Schauspieler}|pw}{}\ledrightnote{\textcolor{blue}{Lucien Guitry}} hat diesen Ursprung – denn ich schrieb die \textcolor{green}{Notiz}\pwindex{Zuckerkandl, Berta 13.\,4.\,1864 Wien – 16.\,10.\,1945 Paris@\textsc{Zuckerkandl, Berta} (13.\,4.\,1864 Wien – 16.\,10.\,1945 Paris), \emph{Schriftstellerin, Journalistin, Übersetzerin}!Courrier des Théatres. Au jour le jour. Shakespeare disait [...] [Über die Wiener Uraufführung von Das weite Land]@\strich\emph{Courrier des Théatres. Au jour le jour. Shakespeare disait [...] [Über die Wiener Uraufführung von Das weite Land]}|pwv}{}\ledrightnote{{$\rightarrow$}\emph{\textcolor{green}{Courrier des Théatres. Au jour le jour. Shakespeare disait [...] [Über die Wiener Uraufführung von Das weite Land]}}} eben so – dass \textcolor{blue}{Guitry}\pwindex{Guitry, Lucien 13.\,12.\,1860 Paris – 1.\,6.\,1925 ebd.@\textsc{Guitry, Lucien} (13.\,12.\,1860 Paris – 1.\,6.\,1925 ebd.), \emph{Schriftsteller, Schauspieler}|pw}{}\ledrightnote{\textcolor{blue}{Lucien Guitry}} aufmerksam werden musste.\pend
           
\pstart
           {\pb}Selbstverständlich möchte ich aber gar
               nicht eine bessere Möglichkeit stören. Und wenn Sie glauben dass Herr \textcolor{blue}{Rémon}\pwindex{Rémon, Maurice 27.\,11.\,1861 Paris – 20.\,6.\,1945 Mérignac@\textsc{Rémon, Maurice} (27.\,11.\,1861 Paris – 20.\,6.\,1945 Mérignac), \emph{Übersetzer}|pw}{}\ledrightnote{\textcolor{blue}{Maurice Rémon}} rascher zum Ziel ko{\geminationm}t – so trete ich gerne zurück. Die Haupt{\pb}sache bleibt die \textcolor{green}{Aufführung}\pwindex{Schnitzler, Arthur 15. 5. 1862 Wien – 21. 10. 1931 ebd.@\textsc{Schnitzler, Arthur} (15. 5. 1862 Wien – 21. 10. 1931 ebd.), \emph{Schriftsteller, Mediziner}!weite Land. Tragikomödie in fünf Akten@\strich\emph{Das weite Land. Tragikomödie in fünf Akten}|pwv}{}\ledrightnote{{$\rightarrow$}\emph{\textcolor{green}{Das weite Land. Tragikomödie in fünf Akten}}}. Sollte aber Herr \textcolor{blue}{Rémon}\pwindex{Rémon, Maurice 27.\,11.\,1861 Paris – 20.\,6.\,1945 Mérignac@\textsc{Rémon, Maurice} (27.\,11.\,1861 Paris – 20.\,6.\,1945 Mérignac), \emph{Übersetzer}|pw}{}\ledrightnote{\textcolor{blue}{Maurice Rémon}} sich mit mir in Verbindung setzen wollen
               – so bin ich auch dazu bereit.\pend
           
\pstart
           – Bitte handeln Sie (ohne jede Behinderung) ganz nach bester Einsicht.\pend
           
\pstart
           Herzlich{[}s{]}t grüssend {\\[\baselineskip]}\spacefill\mbox{Berta Zuckerk}\pend
           \leftskip=0em{}\selectlanguage{ngerman}\vspace{1em}{\vspace{1\baselineskip}}
\pstart
           {\pb}\textcolor{gray}{\textbf{\begin{otherlanguage}{french}\label{K_L04004-4v}\edtext{\textcolor{blue}{Shakespeare}\pwindex{Shakespeare, William 23.\,4.\,1564? Stratford-upon-Avon – 3.\,5.\,1616 ebd.@\textsc{Shakespeare, William} (23.\,4.\,1564? Stratford-upon-Avon – 3.\,5.\,1616 ebd.), \emph{Schauspieler, Dramatiker}|pw}{}\ledrightnote{\textcolor{blue}{William Shakespeare}} disait:}{\lemma{\textnormal{\emph{Shakespeare disait:}}}\Cendnote{\textnormal{französisch: »\textcolor{blue}{Shakespeare}\pwindex{Shakespeare, William 23.\,4.\,1564? Stratford-upon-Avon – 3.\,5.\,1616 ebd.@\textsc{Shakespeare, William} (23.\,4.\,1564? Stratford-upon-Avon – 3.\,5.\,1616 ebd.), \emph{Schauspieler, Dramatiker}|pw} sagte: ›Du weißt, eine fremde Seele, ist
                        ein dunkler Wald{\dots}‹ Für \textcolor{blue}{Arthur Schnitzler}, den \textcolor{pink}{Wiener}\oindex{Wien@\textbf{Wien}, \emph{Verwaltungsgebiet}|pw} Dramatiker (an dessen 
                        \textcolor{green}{\emph{Abschiedssouper}}\pwindex{Schnitzler, Arthur 15. 5. 1862 Wien – 21. 10. 1931 ebd.@\textsc{Schnitzler, Arthur} (15. 5. 1862 Wien – 21. 10. 1931 ebd.), \emph{Schriftsteller, Mediziner}!Abschiedssouper@\strich\emph{Abschiedssouper}|pw}
 man sich erinnert), ist die Seele der Menschen ein ›weites Land‹.
                        Unter diesem Titel hat er gerade am \textcolor{brown}{Burgtheater}\orgindex{Burgtheater@Burgtheater|pw} in \textcolor{pink}{Wien}\oindex{Wien@\textbf{Wien}, \emph{Verwaltungsgebiet}|pw} mit großem
                        Erfolg eine \textcolor{green}{Tragikomödie}\pwindex{Schnitzler, Arthur 15. 5. 1862 Wien – 21. 10. 1931 ebd.@\textsc{Schnitzler, Arthur} (15. 5. 1862 Wien – 21. 10. 1931 ebd.), \emph{Schriftsteller, Mediziner}!weite Land. Tragikomödie in fünf Akten@\strich\emph{Das weite Land. Tragikomödie in fünf Akten}|pwv} gegeben, die ganz \textcolor{pink}{Europa}\oindex{Europa@\textbf{Europa}|pw} erobern wird. Die Psychologie der Figuren und die
                        feinsinnige und tiefgründige Beobachtung des menschlichen Bewusstseins
                        werden diesem \textcolor{green}{Stück}\pwindex{Schnitzler, Arthur 15. 5. 1862 Wien – 21. 10. 1931 ebd.@\textsc{Schnitzler, Arthur} (15. 5. 1862 Wien – 21. 10. 1931 ebd.), \emph{Schriftsteller, Mediziner}!weite Land. Tragikomödie in fünf Akten@\strich\emph{Das weite Land. Tragikomödie in fünf Akten}|pwv},
                        einem Meisterwerk, einen besonderen Platz in der Hochachtung aller Kenner
                        sichern. Die Hauptrolle, ein großer Industrieller, war für \textcolor{blue}{Joseph Kainz}\pwindex{Kainz, Josef 2.\,1.\,1858 Mosonmagyaróvár – 20.\,9.\,1910 Wien@\textsc{Kainz, Josef} (2.\,1.\,1858 Mosonmagyaróvár – 20.\,9.\,1910 Wien), \emph{Schauspieler}|pw} bestimmt, den großen Schauspieler, um
                        dessen Verlust \textcolor{pink}{Österreich}\oindex{Österreich-Ungarn@\textbf{Österreich-Ungarn}|pw} noch immer
                        trauert!«}}}\label{K_L04004-4} »\label{K_L04004-5v}\edtext{L’âme
                     d’autrui, vois-tu, c’est une forêt obscure}{\lemma{\textnormal{\emph{L’âme … obscure}}}\Cendnote{\textnormal{Das Zitat stammt nicht aus \textcolor{blue}{Shakespeares}\pwindex{Shakespeare, William 23.\,4.\,1564? Stratford-upon-Avon – 3.\,5.\,1616 ebd.@\textsc{Shakespeare, William} (23.\,4.\,1564? Stratford-upon-Avon – 3.\,5.\,1616 ebd.), \emph{Schauspieler, Dramatiker}|pwk} Werken sondern aus dem 17. Kapitel von
                           \textcolor{blue}{Ivan Sergeevič Turgenevs}\pwindex{Turgenev, Ivan Sergeevič 9.\,11.\,1818 Orjol – 3.\,9.\,1883 Bougival@\textsc{Turgenev, Ivan Sergeevič} (9.\,11.\,1818 Orjol – 3.\,9.\,1883 Bougival), \emph{Schriftsteller}|pwk} Roman \emph{\textcolor{green}{Dvorjanskoe gnezdo}\pwindex{Turgenev, Ivan Sergeevič 9.\,11.\,1818 Orjol – 3.\,9.\,1883 Bougival@\textsc{Turgenev, Ivan Sergeevič} (9.\,11.\,1818 Orjol – 3.\,9.\,1883 Bougival), \emph{Schriftsteller}!Dvorjanskoe gnezdo@\strich\emph{Dvorjanskoe gnezdo}|pwk}} (deutsch: \emph{\textcolor{green}{Das adelige Nest}\pwindex{Turgenev, Ivan Sergeevič 9.\,11.\,1818 Orjol – 3.\,9.\,1883 Bougival@\textsc{Turgenev, Ivan Sergeevič} (9.\,11.\,1818 Orjol – 3.\,9.\,1883 Bougival), \emph{Schriftsteller}!adelige Nest@\strich\emph{Das adelige Nest}|pwk}}).}}}\label{K_L04004-5}{\dots}« Pour M. \textcolor{blue}{Arthur
                        Schnitzler}{}\ledrightnote{}, l’auteur dramatique \textcolor{pink}{viennois}\oindex{Wien@\textbf{Wien}, \emph{Verwaltungsgebiet}|pw}{}\ledrightnote{\textcolor{pink}{Wien}} (dont on se rappelle \emph{\textcolor{green}{Souper d’adieu}\pwindex{Schnitzler, Arthur 15. 5. 1862 Wien – 21. 10. 1931 ebd.@\textsc{Schnitzler, Arthur} (15. 5. 1862 Wien – 21. 10. 1931 ebd.), \emph{Schriftsteller, Mediziner}!Souper d’Adieu@\strich\emph{Souper d’Adieu}|pw}\pwindex{Schnitzler, Arthur 15. 5. 1862 Wien – 21. 10. 1931 ebd.@\textsc{Schnitzler, Arthur} (15. 5. 1862 Wien – 21. 10. 1931 ebd.), \emph{Schriftsteller, Mediziner}!Abschiedssouper@\strich\emph{Abschiedssouper}|pw}{}\ledrightnote{\textcolor{green}{Souper d’Adieu}{\newline}\textcolor{green}{Abschiedssouper}}}), l’âme des hommes est un »vaste pays«. Sous ce titre, il vient de donner
                     au \textcolor{brown}{théâtre de la Cour}\orgindex{Burgtheater@Burgtheater|pw}{}\ledrightnote{\textcolor{brown}{Burgtheater}}, à \textcolor{pink}{Vienne}\oindex{Wien@\textbf{Wien}, \emph{Verwaltungsgebiet}|pw}{}\ledrightnote{\textcolor{pink}{Wien}}, avec un succès retentissant, une \textcolor{green}{tragi-comédie}\pwindex{Schnitzler, Arthur 15. 5. 1862 Wien – 21. 10. 1931 ebd.@\textsc{Schnitzler, Arthur} (15. 5. 1862 Wien – 21. 10. 1931 ebd.), \emph{Schriftsteller, Mediziner}!weite Land. Tragikomödie in fünf Akten@\strich\emph{Das weite Land. Tragikomödie in fünf Akten}|pwv}{}\ledrightnote{{$\rightarrow$}\emph{\textcolor{green}{Das weite Land. Tragikomödie in fünf Akten}}} qui fera le tour de
                        l’\textcolor{pink}{Europe}\oindex{Europa@\textbf{Europa}|pw}{}\ledrightnote{\textcolor{pink}{Europa}}. La psychologie des caractères,
                     l’observation qui s’y remarque, délicate et profonde, de la conscience humaine,
                     assureront à cette \textcolor{green}{pièce}\pwindex{Schnitzler, Arthur 15. 5. 1862 Wien – 21. 10. 1931 ebd.@\textsc{Schnitzler, Arthur} (15. 5. 1862 Wien – 21. 10. 1931 ebd.), \emph{Schriftsteller, Mediziner}!weite Land. Tragikomödie in fünf Akten@\strich\emph{Das weite Land. Tragikomödie in fünf Akten}|pwv}{}\ledrightnote{{$\rightarrow$}\emph{\textcolor{green}{Das weite Land. Tragikomödie in fünf Akten}}}, qui est une œuvre, une place à part dans l’estime de tous les
                     connaisseurs. Le rôle principal, un type de grand industriel, était destiné à
                        \textcolor{blue}{Joseph \label{T_L04004-1v}\edtext{Kaïnz}{\lemma{\textnormal{\emph{Kaïnz}}}\Cendnote{\textnormal{Im
                           Text steht an dieser und der folgenden Erwähnung: »Haïnz«.}}}\label{T_L04004-1}}\pwindex{Kainz, Josef 2.\,1.\,1858 Mosonmagyaróvár – 20.\,9.\,1910 Wien@\textsc{Kainz, Josef} (2.\,1.\,1858 Mosonmagyaróvár – 20.\,9.\,1910 Wien), \emph{Schauspieler}|pw}{}\ledrightnote{\textcolor{blue}{Josef Kainz}}, le grand comédien dont l’\textcolor{pink}{Autriche}\oindex{Österreich-Ungarn@\textbf{Österreich-Ungarn}|pw}{}\ledrightnote{\textcolor{pink}{Österreich-Ungarn}}
                     pleure encore la perte!\end{otherlanguage}}}\pend
           
\pstart
           \textcolor{gray}{\textbf{\begin{otherlanguage}{french}– 
                     \label{K_L04004-6v}\edtext{Enfin, voilà donc une \textcolor{green}{pièce}\pwindex{Schnitzler, Arthur 15. 5. 1862 Wien – 21. 10. 1931 ebd.@\textsc{Schnitzler, Arthur} (15. 5. 1862 Wien – 21. 10. 1931 ebd.), \emph{Schriftsteller, Mediziner}!weite Land. Tragikomödie in fünf Akten@\strich\emph{Das weite Land. Tragikomödie in fünf Akten}|pwv}{}\ledrightnote{{$\rightarrow$}\emph{\textcolor{green}{Das weite Land. Tragikomödie in fünf Akten}}} ; une création
                     vraiment belle à tenter! s’écria ce dernier sur son lit de mort, quand on lui
                     lut \emph{\textcolor{green}{le Vaste pays}\pwindex{Schnitzler, Arthur 15. 5. 1862 Wien – 21. 10. 1931 ebd.@\textsc{Schnitzler, Arthur} (15. 5. 1862 Wien – 21. 10. 1931 ebd.), \emph{Schriftsteller, Mediziner}!weite Land. Tragikomödie in fünf Akten@\strich\emph{Das weite Land. Tragikomödie in fünf Akten}|pw}{}\ledrightnote{\textcolor{green}{Das weite Land. Tragikomödie in fünf Akten}}}}{\lemma{\textnormal{\emph{Enfin, … pays}}}\Cendnote{\textnormal{französisch: »Endlich, hier ist
                        also ein \textcolor{green}{Stück}\pwindex{Schnitzler, Arthur 15. 5. 1862 Wien – 21. 10. 1931 ebd.@\textsc{Schnitzler, Arthur} (15. 5. 1862 Wien – 21. 10. 1931 ebd.), \emph{Schriftsteller, Mediziner}!weite Land. Tragikomödie in fünf Akten@\strich\emph{Das weite Land. Tragikomödie in fünf Akten}|pwv}, ein
                        wirklich schönes Werk, das einen Versuch wert ist! rief er auf seinem
                        Sterbebett aus, als man ihm \emph{\textcolor{green}{Das weite Land}\pwindex{Schnitzler, Arthur 15. 5. 1862 Wien – 21. 10. 1931 ebd.@\textsc{Schnitzler, Arthur} (15. 5. 1862 Wien – 21. 10. 1931 ebd.), \emph{Schriftsteller, Mediziner}!weite Land. Tragikomödie in fünf Akten@\strich\emph{Das weite Land. Tragikomödie in fünf Akten}|pw}} vorlas.«}}}\label{K_L04004-6}.\end{otherlanguage}}}\pend
           
\pstart
           \textcolor{gray}{\textbf{\begin{otherlanguage}{french}\label{K_L04004-7v}\edtext{\textcolor{blue}{Joseph Kaïnz}\pwindex{Kainz, Josef 2.\,1.\,1858 Mosonmagyaróvár – 20.\,9.\,1910 Wien@\textsc{Kainz, Josef} (2.\,1.\,1858 Mosonmagyaróvár – 20.\,9.\,1910 Wien), \emph{Schauspieler}|pw}{}\ledrightnote{\textcolor{blue}{Josef Kainz}} eût été superbe dans \emph{\textcolor{green}{le Vaste Pays}\pwindex{Schnitzler, Arthur 15. 5. 1862 Wien – 21. 10. 1931 ebd.@\textsc{Schnitzler, Arthur} (15. 5. 1862 Wien – 21. 10. 1931 ebd.), \emph{Schriftsteller, Mediziner}!weite Land. Tragikomödie in fünf Akten@\strich\emph{Das weite Land. Tragikomödie in fünf Akten}|pw}{}\ledrightnote{\textcolor{green}{Das weite Land. Tragikomödie in fünf Akten}}}. Il eût mis dans tout leur relief les côtés hautains, douloureux, amers,
                     passionnés du personnage. Pour jouer en \textcolor{pink}{France}\oindex{Frankreich@\textbf{Frankreich}|pw}{}\ledrightnote{\textcolor{pink}{Frankreich}} et en \textcolor{pink}{Europe}\oindex{Europa@\textbf{Europa}|pw}{}\ledrightnote{\textcolor{pink}{Europa}} un tel rôle,
                     il faudrait \textcolor{blue}{Lucien Guitry}\pwindex{Guitry, Lucien 13.\,12.\,1860 Paris – 1.\,6.\,1925 ebd.@\textsc{Guitry, Lucien} (13.\,12.\,1860 Paris – 1.\,6.\,1925 ebd.), \emph{Schriftsteller, Schauspieler}|pw}{}\ledrightnote{\textcolor{blue}{Lucien Guitry}}; lui seul
                     aurait la puissance de talent nécessaire à cette création. Pendant les
                     entr’actes de la \textcolor{green}{pièce}\pwindex{Schnitzler, Arthur 15. 5. 1862 Wien – 21. 10. 1931 ebd.@\textsc{Schnitzler, Arthur} (15. 5. 1862 Wien – 21. 10. 1931 ebd.), \emph{Schriftsteller, Mediziner}!weite Land. Tragikomödie in fünf Akten@\strich\emph{Das weite Land. Tragikomödie in fünf Akten}|pw}{}\ledrightnote{\textcolor{green}{Das weite Land. Tragikomödie in fünf Akten}}, le
                        soir de la \textcolor{violet}{première}\eventindex{Burgtheater@\textbf{Burgtheater}!Premiere von Das weite Land, 14.10.1911 [I.]@Premiere von Das weite Land, 14.10.1911 [I.]|pwv}{}\ledrightnote{{$\rightarrow$}\emph{\textcolor{violet}{Premiere von Das weite Land, 14.10.1911 [I.]}}}, il n’ y avait qu’une voix : »Voilà un rôle pour \textcolor{blue}{Lucien Guitry}\pwindex{Guitry, Lucien 13.\,12.\,1860 Paris – 1.\,6.\,1925 ebd.@\textsc{Guitry, Lucien} (13.\,12.\,1860 Paris – 1.\,6.\,1925 ebd.), \emph{Schriftsteller, Schauspieler}|pw}{}\ledrightnote{\textcolor{blue}{Lucien Guitry}}! Comme \textcolor{blue}{Guitry}\pwindex{Guitry, Lucien 13.\,12.\,1860 Paris – 1.\,6.\,1925 ebd.@\textsc{Guitry, Lucien} (13.\,12.\,1860 Paris – 1.\,6.\,1925 ebd.), \emph{Schriftsteller, Schauspieler}|pw}{}\ledrightnote{\textcolor{blue}{Lucien Guitry}} serait beau dans ce personnage!{\dots}« Hommage précieux, hommage significatif rendu au
                     grand \textcolor{blue}{artiste}\pwindex{Guitry, Lucien 13.\,12.\,1860 Paris – 1.\,6.\,1925 ebd.@\textsc{Guitry, Lucien} (13.\,12.\,1860 Paris – 1.\,6.\,1925 ebd.), \emph{Schriftsteller, Schauspieler}|pwv}{}\ledrightnote{{$\rightarrow$}\emph{\textcolor{blue}{Lucien Guitry}}} que
                        l’\textcolor{pink}{Europe}\oindex{Europa@\textbf{Europa}|pw}{}\ledrightnote{\textcolor{pink}{Europa}} entière nous envie!}{\lemma{\textnormal{\emph{Joseph … envie!}}}\Cendnote{\textnormal{französisch: »\textcolor{blue}{Joseph Kainz}\pwindex{Kainz, Josef 2.\,1.\,1858 Mosonmagyaróvár – 20.\,9.\,1910 Wien@\textsc{Kainz, Josef} (2.\,1.\,1858 Mosonmagyaróvár – 20.\,9.\,1910 Wien), \emph{Schauspieler}|pw} wäre großartig in \emph{\textcolor{green}{Das weite Land}\pwindex{Schnitzler, Arthur 15. 5. 1862 Wien – 21. 10. 1931 ebd.@\textsc{Schnitzler, Arthur} (15. 5. 1862 Wien – 21. 10. 1931 ebd.), \emph{Schriftsteller, Mediziner}!weite Land. Tragikomödie in fünf Akten@\strich\emph{Das weite Land. Tragikomödie in fünf Akten}|pw}} gewesen. Er hätte die hochmütigen, schmerzhaften, die bitteren und
                        leidenschaftlichen Seiten der Figur voll zur Geltung gebracht. Um eine
                        solche Rolle in \textcolor{pink}{Frankreich}\oindex{Frankreich@\textbf{Frankreich}|pw} und \textcolor{pink}{Europa}\oindex{Europa@\textbf{Europa}|pw} zu spielen, bräuchte man \textcolor{blue}{Lucien Guitry}\pwindex{Guitry, Lucien 13.\,12.\,1860 Paris – 1.\,6.\,1925 ebd.@\textsc{Guitry, Lucien} (13.\,12.\,1860 Paris – 1.\,6.\,1925 ebd.), \emph{Schriftsteller, Schauspieler}|pw}; nur er allein hätte das
                        nötige Talent für diese Kunst. Während der Pausen des \textcolor{green}{Stücks}\pwindex{Schnitzler, Arthur 15. 5. 1862 Wien – 21. 10. 1931 ebd.@\textsc{Schnitzler, Arthur} (15. 5. 1862 Wien – 21. 10. 1931 ebd.), \emph{Schriftsteller, Mediziner}!weite Land. Tragikomödie in fünf Akten@\strich\emph{Das weite Land. Tragikomödie in fünf Akten}|pw} am Abend der \textcolor{violet}{Premiere}\eventindex{Burgtheater@\textbf{Burgtheater}!Premiere von Das weite Land, 14.10.1911 [I.]@Premiere von Das weite Land, 14.10.1911 [I.]|pwv} gab es nur eine Stimme:
                        ›Das ist eine Rolle für \textcolor{blue}{Lucien
                        Guitry}\pwindex{Guitry, Lucien 13.\,12.\,1860 Paris – 1.\,6.\,1925 ebd.@\textsc{Guitry, Lucien} (13.\,12.\,1860 Paris – 1.\,6.\,1925 ebd.), \emph{Schriftsteller, Schauspieler}|pw}! Wie gut \textcolor{blue}{Guitry}\pwindex{Guitry, Lucien 13.\,12.\,1860 Paris – 1.\,6.\,1925 ebd.@\textsc{Guitry, Lucien} (13.\,12.\,1860 Paris – 1.\,6.\,1925 ebd.), \emph{Schriftsteller, Schauspieler}|pw} in dieser
                        Rolle wäre! {\dots}‹ Eine prachtvolle Hommage, eine
                        aussagekräftige Hommage an den großen \textcolor{blue}{Künstler}\pwindex{Guitry, Lucien 13.\,12.\,1860 Paris – 1.\,6.\,1925 ebd.@\textsc{Guitry, Lucien} (13.\,12.\,1860 Paris – 1.\,6.\,1925 ebd.), \emph{Schriftsteller, Schauspieler}|pwv}, um den uns ganz \textcolor{pink}{Europa}\oindex{Europa@\textbf{Europa}|pw} beneidet!«}}}\label{K_L04004-7}\end{otherlanguage}}}\pend
           \selectlanguage{ngerman}\endnumbering\briefempfaengerindex{Schnitzler, Arthur@\textsc{Schnitzler, Arthur}!zzzZuckerkandl, Berta@\emph{von Berta Zuckerkandl}!1911-11-072@{{[}7. 11. 1911?{]}}|)be}\mylabel{L04004h}  \newcommand{\dateiname}{L04004}\newcommand{\titel}{Berta Zuckerkandl an Arthur Schnitzler, [7. 11. 1911?]}\newcommand{\editorInnen}{Herausgegeben von Jahnke, SelmaMüller, Martin Anton}\input{../tex-inputs/latex-korrekturansicht-abspann}
